% Options for packages loaded elsewhere
\PassOptionsToPackage{unicode}{hyperref}
\PassOptionsToPackage{hyphens}{url}
\documentclass[
]{ltjsarticle}
\usepackage{xcolor}
\usepackage[margin=2cm]{geometry}
\usepackage{amsmath,amssymb}
\setcounter{secnumdepth}{-\maxdimen} % remove section numbering
\usepackage{iftex}
\ifPDFTeX
  \usepackage[T1]{fontenc}
  \usepackage[utf8]{inputenc}
  \usepackage{textcomp} % provide euro and other symbols
\else % if luatex or xetex
  \usepackage{unicode-math} % this also loads fontspec
  \defaultfontfeatures{Scale=MatchLowercase}
  \defaultfontfeatures[\rmfamily]{Ligatures=TeX,Scale=1}
\fi
\usepackage{lmodern}
\ifPDFTeX\else
  % xetex/luatex font selection
  \setmainfont[]{Hiragino Sans}
  \setmonofont[]{Menlo}
\fi
% Use upquote if available, for straight quotes in verbatim environments
\IfFileExists{upquote.sty}{\usepackage{upquote}}{}
\IfFileExists{microtype.sty}{% use microtype if available
  \usepackage[]{microtype}
  \UseMicrotypeSet[protrusion]{basicmath} % disable protrusion for tt fonts
}{}
\makeatletter
\@ifundefined{KOMAClassName}{% if non-KOMA class
  \IfFileExists{parskip.sty}{%
    \usepackage{parskip}
  }{% else
    \setlength{\parindent}{0pt}
    \setlength{\parskip}{6pt plus 2pt minus 1pt}}
}{% if KOMA class
  \KOMAoptions{parskip=half}}
\makeatother
\usepackage{color}
\usepackage{fancyvrb}
\newcommand{\VerbBar}{|}
\newcommand{\VERB}{\Verb[commandchars=\\\{\}]}
\DefineVerbatimEnvironment{Highlighting}{Verbatim}{commandchars=\\\{\}}
% Add ',fontsize=\small' for more characters per line
\newenvironment{Shaded}{}{}
\newcommand{\AlertTok}[1]{\textcolor[rgb]{1.00,0.00,0.00}{\textbf{#1}}}
\newcommand{\AnnotationTok}[1]{\textcolor[rgb]{0.38,0.63,0.69}{\textbf{\textit{#1}}}}
\newcommand{\AttributeTok}[1]{\textcolor[rgb]{0.49,0.56,0.16}{#1}}
\newcommand{\BaseNTok}[1]{\textcolor[rgb]{0.25,0.63,0.44}{#1}}
\newcommand{\BuiltInTok}[1]{\textcolor[rgb]{0.00,0.50,0.00}{#1}}
\newcommand{\CharTok}[1]{\textcolor[rgb]{0.25,0.44,0.63}{#1}}
\newcommand{\CommentTok}[1]{\textcolor[rgb]{0.38,0.63,0.69}{\textit{#1}}}
\newcommand{\CommentVarTok}[1]{\textcolor[rgb]{0.38,0.63,0.69}{\textbf{\textit{#1}}}}
\newcommand{\ConstantTok}[1]{\textcolor[rgb]{0.53,0.00,0.00}{#1}}
\newcommand{\ControlFlowTok}[1]{\textcolor[rgb]{0.00,0.44,0.13}{\textbf{#1}}}
\newcommand{\DataTypeTok}[1]{\textcolor[rgb]{0.56,0.13,0.00}{#1}}
\newcommand{\DecValTok}[1]{\textcolor[rgb]{0.25,0.63,0.44}{#1}}
\newcommand{\DocumentationTok}[1]{\textcolor[rgb]{0.73,0.13,0.13}{\textit{#1}}}
\newcommand{\ErrorTok}[1]{\textcolor[rgb]{1.00,0.00,0.00}{\textbf{#1}}}
\newcommand{\ExtensionTok}[1]{#1}
\newcommand{\FloatTok}[1]{\textcolor[rgb]{0.25,0.63,0.44}{#1}}
\newcommand{\FunctionTok}[1]{\textcolor[rgb]{0.02,0.16,0.49}{#1}}
\newcommand{\ImportTok}[1]{\textcolor[rgb]{0.00,0.50,0.00}{\textbf{#1}}}
\newcommand{\InformationTok}[1]{\textcolor[rgb]{0.38,0.63,0.69}{\textbf{\textit{#1}}}}
\newcommand{\KeywordTok}[1]{\textcolor[rgb]{0.00,0.44,0.13}{\textbf{#1}}}
\newcommand{\NormalTok}[1]{#1}
\newcommand{\OperatorTok}[1]{\textcolor[rgb]{0.40,0.40,0.40}{#1}}
\newcommand{\OtherTok}[1]{\textcolor[rgb]{0.00,0.44,0.13}{#1}}
\newcommand{\PreprocessorTok}[1]{\textcolor[rgb]{0.74,0.48,0.00}{#1}}
\newcommand{\RegionMarkerTok}[1]{#1}
\newcommand{\SpecialCharTok}[1]{\textcolor[rgb]{0.25,0.44,0.63}{#1}}
\newcommand{\SpecialStringTok}[1]{\textcolor[rgb]{0.73,0.40,0.53}{#1}}
\newcommand{\StringTok}[1]{\textcolor[rgb]{0.25,0.44,0.63}{#1}}
\newcommand{\VariableTok}[1]{\textcolor[rgb]{0.10,0.09,0.49}{#1}}
\newcommand{\VerbatimStringTok}[1]{\textcolor[rgb]{0.25,0.44,0.63}{#1}}
\newcommand{\WarningTok}[1]{\textcolor[rgb]{0.38,0.63,0.69}{\textbf{\textit{#1}}}}
\usepackage{longtable,booktabs,array}
\newcounter{none} % for unnumbered tables
\usepackage{calc} % for calculating minipage widths
% Correct order of tables after \paragraph or \subparagraph
\usepackage{etoolbox}
\makeatletter
\patchcmd\longtable{\par}{\if@noskipsec\mbox{}\fi\par}{}{}
\makeatother
% Allow footnotes in longtable head/foot
\IfFileExists{footnotehyper.sty}{\usepackage{footnotehyper}}{\usepackage{footnote}}
\makesavenoteenv{longtable}
\setlength{\emergencystretch}{3em} % prevent overfull lines
\providecommand{\tightlist}{%
  \setlength{\itemsep}{0pt}\setlength{\parskip}{0pt}}
\usepackage{bookmark}
\IfFileExists{xurl.sty}{\usepackage{xurl}}{} % add URL line breaks if available
\urlstyle{same}
\hypersetup{
  hidelinks,
  pdfcreator={LaTeX via pandoc}}

\author{}
\date{}

\begin{document}

{
\setcounter{tocdepth}{2}
\tableofcontents
}
\section{PsiLang --- Minimal ML-like Functional Language for
Algorithms}\label{psilang-minimal-ml-like-functional-language-for-algorithms}

進化計算 (本ディレクトリ親 \texttt{../REPORT.md} 参照) の top-5
個体から抽出した\textbf{コンセンサス遺伝子}を実装したもの。v1 の PhiLang
が orchestration shell (\texttt{compute\ =\ X} の 1 行 dispatch)
だったのに対し、PsiLang
は\textbf{アルゴリズムを言語自身の構文で書ける}汎用 mini-ML である。

\subsection{構成}\label{ux69cbux6210}

{\def\LTcaptype{none} % do not increment counter
\begin{longtable}[]{@{}
  >{\raggedright\arraybackslash}p{(\linewidth - 4\tabcolsep) * \real{0.3000}}
  >{\raggedright\arraybackslash}p{(\linewidth - 4\tabcolsep) * \real{0.3000}}
  >{\raggedleft\arraybackslash}p{(\linewidth - 4\tabcolsep) * \real{0.4000}}@{}}
\toprule\noalign{}
\begin{minipage}[b]{\linewidth}\raggedright
ファイル
\end{minipage} & \begin{minipage}[b]{\linewidth}\raggedright
役割
\end{minipage} & \begin{minipage}[b]{\linewidth}\raggedleft
行数
\end{minipage} \\
\midrule\noalign{}
\endhead
\bottomrule\noalign{}
\endlastfoot
\texttt{psilang.py} & Lexer + Parser + AST + Tree-walker Evaluator +
Stdlib (単一ファイル) & \textasciitilde570 \\
\texttt{test\_basic.psi} & 動作確認: factorial / list / pattern match /
Rat & 15 \\
\texttt{chudnovsky.psi} & \textbf{π 10,000 桁を binary splitting で計算}
& 50 \\
\texttt{pi\_10k.txt} & 実行結果 (3 + 10,000 digits) & --- \\
\texttt{README.md} & この文書 & --- \\
\end{longtable}
}

\subsection{言語仕様
(概要)}\label{ux8a00ux8a9eux4ed5ux69d8-ux6982ux8981}

\begin{itemize}
\tightlist
\item
  \textbf{paradigm:} functional (純粋関数ではない --- 副作用は
  \texttt{emit}/\texttt{print} に限定)
\item
  \textbf{type\_system:} static\_simple
  (型注釈は構文上受け取るが実行時は無視)
\item
  \textbf{arithmetic\_primitives:} \texttt{Int} (任意精度) +
  \texttt{Rat} (有理数、自動約分) + \texttt{Real}
  (任意精度小数、\texttt{Decimal} 由来)
\item
  \textbf{control\_flow:} \texttt{if} / \texttt{match} /
  general\_recursion (no loops; iteration はパターンマッチ + 再帰)
\item
  \textbf{syntax\_family:} ml\_like (let / fn / match / \textbar)
\end{itemize}

\subsubsection{文法 (要約)}\label{ux6587ux6cd5-ux8981ux7d04}

\begin{verbatim}
fn name (p1: T1, p2: T2) -> Ret = body

let x = expr in body
let (a, b) = expr in body

if cond then a else b
match x with | pat1 -> e1 | pat2 -> e2 end

(a, b)         # 2-tuple
[1, 2, 3]      # list
x :: xs        # list cons, []  empty list
\end{verbatim}

\subsubsection{Builtin (一部)}\label{builtin-ux4e00ux90e8}

\begin{itemize}
\tightlist
\item
  \textbf{Int:} \texttt{+\ -\ *\ /\ mod\ ==}, \texttt{int\_fact},
  \texttt{int\_pow}, \texttt{int\_abs}, \texttt{int\_sign},
  \texttt{int\_to\_string}
\item
  \textbf{Rat:} \texttt{rat(n,\ d)}, \texttt{rat\_add},
  \texttt{rat\_sub}, \texttt{rat\_mul}, \texttt{rat\_div},
  \texttt{rat\_neg}, \texttt{rat\_inv}, \texttt{rat\_num},
  \texttt{rat\_den}, \texttt{rat\_from\_int}
\item
  \textbf{Real:} \texttt{real\_from\_int},
  \texttt{real\_from\_rat(r,\ prec)}, \texttt{real\_sqrt(r,\ prec)},
  \texttt{real\_mul}, \texttt{real\_div},
  \texttt{real\_to\_string(r,\ n)}, \texttt{set\_precision}
\item
  \textbf{Misc:} \texttt{fst}, \texttt{snd}, \texttt{length},
  \texttt{emit}, \texttt{print}
\end{itemize}

\subsection{動作確認}\label{ux52d5ux4f5cux78baux8a8d}

\begin{Shaded}
\begin{Highlighting}[]
\ExtensionTok{/usr/bin/python3}\NormalTok{ psilang.py test\_basic.psi}
\CommentTok{\# fact(20) = 2432902008176640000}
\CommentTok{\# sum([1..5]) = 15}
\CommentTok{\# 1+1/2+1/3+1/4 = 25/12}
\CommentTok{\# int\_fact(15) = 1307674368000}
\CommentTok{\# int\_pow(2, 10) = 1024}
\end{Highlighting}
\end{Shaded}

\subsection{π 10,000
桁の計算}\label{ux3c0-10000-ux6841ux306eux8a08ux7b97}

\begin{Shaded}
\begin{Highlighting}[]
\ExtensionTok{/usr/bin/python3}\NormalTok{ psilang.py chudnovsky.psi }\OperatorTok{\textgreater{}}\NormalTok{ pi\_10k.txt}
\CommentTok{\# 実時間: 3.03 秒 (Apple M2)}
\end{Highlighting}
\end{Shaded}

\textbf{\texttt{chudnovsky.psi} の中身} (50 行) は
\textbf{アルゴリズムそのもの} を PsiLang で書いている:

\begin{verbatim}
fn term(k) =
  let sign = if k mod 2 == 0 then 1 else -1 in
  let num  = sign * int_fact(6 * k) * (545140134 * k + 13591409) in
  let den  = int_fact(3 * k) * int_pow(int_fact(k), 3) * int_pow(640320, 3 * k) in
  rat(num, den)

fn bsplit(a, b) =
  if b - a == 1 then term(a)
  else
    let m = (a + b) / 2 in
    let left  = bsplit(a, m) in
    let right = bsplit(m, b) in
    rat_add(left, right)

fn chudnovsky_pi(n_digits) =
  let n_terms = n_digits / 14 + 2 in
  let prec    = n_digits + 20 in
  let _bump   = set_precision(prec) in
  let s_rat   = bsplit(0, n_terms) in
  let s_real  = real_from_rat(s_rat, prec) in
  let c       = real_mul(real_from_int(426880),
                          real_sqrt(real_from_int(10005), prec)) in
  real_div(c, s_real)

let pi  = chudnovsky_pi(10000) in
emit(real_to_string(pi, 10000))
\end{verbatim}

v1 の PhiLang のように
\texttt{compute\ =\ chudnovsky\_binary\_splitting} で逃げず、

\begin{itemize}
\tightlist
\item
  二分木再帰 (\texttt{bsplit})
\item
  Chudnovsky 級数の k 番目項 (\texttt{term})
\item
  最終的な π 復元 (426880 · √10005 / S)
\end{itemize}

をすべて \textbf{PsiLang のソースコード自身が書いている}。実行は
tree-walker interpreter (\texttt{psilang.py}) が Python の
bigint・\texttt{Decimal} を上から借りるだけ。

\subsection{検証結果
(2026-05-16)}\label{ux691cux8a3cux7d50ux679c-2026-05-16}

{\def\LTcaptype{none} % do not increment counter
\begin{longtable}[]{@{}
  >{\raggedright\arraybackslash}p{(\linewidth - 2\tabcolsep) * \real{0.5000}}
  >{\raggedright\arraybackslash}p{(\linewidth - 2\tabcolsep) * \real{0.5000}}@{}}
\toprule\noalign{}
\begin{minipage}[b]{\linewidth}\raggedright
検査項目
\end{minipage} & \begin{minipage}[b]{\linewidth}\raggedright
結果
\end{minipage} \\
\midrule\noalign{}
\endhead
\bottomrule\noalign{}
\endlastfoot
出力サイズ & 10,002 chars = \texttt{"3."} + 10,000 digits \\
先頭 102 桁の既知 π との一致 & ✓ \\
Feynman point (\texttt{999999} の最初の出現位置) & 762 (well-known
値と一致) \\
1000 桁目 & \texttt{9} (OEIS A000796 と一致) \\
v1 PhiLang/\texttt{chudnovsky\_pi(10000)} との bitwise 一致 & ✓ \\
計算時間 & 3.03 秒 (Apple M2) \\
Python interpreter 行数 & 約 650 行 (R4 想定 1500 行を下回る) \\
\end{longtable}
}

\subsection{3 手法の比較 (遅 / 中 /
速)}\label{ux624bux6cd5ux306eux6bd4ux8f03-ux9045-ux4e2d-ux901f}

PsiLang のソースコード\textbf{だけ}で π 10,000 桁を 3 通り (+1 補助)
に計算した時間:

{\def\LTcaptype{none} % do not increment counter
\begin{longtable}[]{@{}
  >{\raggedright\arraybackslash}p{(\linewidth - 10\tabcolsep) * \real{0.1500}}
  >{\raggedright\arraybackslash}p{(\linewidth - 10\tabcolsep) * \real{0.1500}}
  >{\raggedright\arraybackslash}p{(\linewidth - 10\tabcolsep) * \real{0.1500}}
  >{\raggedright\arraybackslash}p{(\linewidth - 10\tabcolsep) * \real{0.1500}}
  >{\raggedleft\arraybackslash}p{(\linewidth - 10\tabcolsep) * \real{0.2000}}
  >{\raggedleft\arraybackslash}p{(\linewidth - 10\tabcolsep) * \real{0.2000}}@{}}
\toprule\noalign{}
\begin{minipage}[b]{\linewidth}\raggedright
\texttt{.psi} ファイル
\end{minipage} & \begin{minipage}[b]{\linewidth}\raggedright
級数
\end{minipage} & \begin{minipage}[b]{\linewidth}\raggedright
内部表現
\end{minipage} & \begin{minipage}[b]{\linewidth}\raggedright
累積
\end{minipage} & \begin{minipage}[b]{\linewidth}\raggedleft
\textbf{実時間}
\end{minipage} & \begin{minipage}[b]{\linewidth}\raggedleft
比率
\end{minipage} \\
\midrule\noalign{}
\endhead
\bottomrule\noalign{}
\endlastfoot
\texttt{machin\_rat.psi} & Machin (1.4 桁/項) & \textbf{Rat} & 線形 &
\textbf{28.24 s} & 100× \\
\texttt{chudnovsky\_linear.psi} & Chudnovsky (14 桁/項) & Rat & 線形 &
3.49 s & 12× \\
\texttt{chudnovsky.psi} & Chudnovsky & Rat & \textbf{二分木} & 3.04 s &
10× \\
\texttt{machin.psi} & Machin & \textbf{Real (Decimal)} & 線形 &
\textbf{0.30 s} & 1× \\
\end{longtable}
}

4 出力すべて \textbf{SHA256 が完全一致}
(\texttt{d44e2dba39a378…})、つまり 10,000 桁すべて bit-wise
に同じ値。詳細は親 \texttt{../REPORT.md} §7.3 を参照。

主な知見: - \textbf{同じ Machin 公式}を Rat で書くと 28 秒、Real
で書くと 0.3 秒 --- 94 倍の差。 - 級数の収束速度より、\textbf{型の選択
(Rat vs Real) が性能を支配}することが PsiLang で可視化される。 -
いずれの \texttt{.psi} も 50 行以下で書ける (R1\_AlgorithmExpressivity
の閾値内)。

\subsection{v1 PhiLang
との対比}\label{v1-philang-ux3068ux306eux5bfeux6bd4}

{\def\LTcaptype{none} % do not increment counter
\begin{longtable}[]{@{}
  >{\raggedright\arraybackslash}p{(\linewidth - 4\tabcolsep) * \real{0.3333}}
  >{\raggedright\arraybackslash}p{(\linewidth - 4\tabcolsep) * \real{0.3333}}
  >{\raggedright\arraybackslash}p{(\linewidth - 4\tabcolsep) * \real{0.3333}}@{}}
\toprule\noalign{}
\begin{minipage}[b]{\linewidth}\raggedright
観点
\end{minipage} & \begin{minipage}[b]{\linewidth}\raggedright
v1 PhiLang
\end{minipage} & \begin{minipage}[b]{\linewidth}\raggedright
v2 PsiLang
\end{minipage} \\
\midrule\noalign{}
\endhead
\bottomrule\noalign{}
\endlastfoot
役割 & Phase orchestration (timeout/retry/checkpoint) & Algorithm
description (recursion/pattern/arithmetic) \\
π 計算ソース & \texttt{compute\ =\ chudnovsky\_binary\_splitting} (1 行
dispatch) & \texttt{chudnovsky.psi} (50 行の真のアルゴリズム) \\
ホスト依存 & Python kernel \texttt{philang\_stdlib.chudnovsky\_pi()}
がすべての仕事 & Python は \texttt{int}, \texttt{Decimal}, gcd
を提供するだけ。algorithm は PsiLang 内 \\
パラダイム & YAML-ish DSL & ML 系関数型 \\
評価モデル & 4 ブロック
(\texttt{phase\ /\ every\ /\ after\ /\ guarantees}) & tree-walker + 再帰
+ パターンマッチ \\
進化計算の根拠 & top 5 共通の 4 つ巴 (bounded burst × mandatory retry ×
\ldots) & top 5 コンセンサス (functional × static\_simple × bigint+rat ×
pattern\_match\_fold × ml\_like) \\
\end{longtable}
}

\subsection{限界}\label{ux9650ux754c}

\begin{itemize}
\tightlist
\item
  \textbf{Closures はサポートしない}: 局所 \texttt{let}
  で定義した関数は環境を捕捉しない (\texttt{fn}
  宣言はトップレベルのみ)。Chudnovsky には不要なので問題なし。
\item
  \textbf{型検査なし}: 型注釈は構文上受け取るが実行時は無視
  (\texttt{static\_simple} を「型注釈は人間用ドキュメント」と解釈)。
\item
  \textbf{末尾呼び出し最適化なし}: 深い線形再帰 (例:
  \texttt{factorial(4000)} を PsiLang 内で純粋再帰で書く) は Python 側で
  stack overflow になる。回避策として \texttt{int\_fact} を builtin
  に置いた。
\item
  \textbf{lambdas (匿名関数) なし}: \texttt{fn}
  は宣言文のみ。FFT/quicksort を書くなら追加が必要。
\item
  \textbf{モジュール / import なし}: 単一ファイルプログラムのみ。
\end{itemize}

\end{document}
